% EGA V, English translation by P. Blass, 1987 (cote 24): archive cote 24, site page 19 (PDF page 20).
% typescript. Draft compared: none.
% AI-transcribed by Claude Opus 5.5 (claude-opus-5-5, adaptive thinking, effort high), Message Batches API, batch ega24b (msgbatch_01RGJJeHdR96U16ZWFuUjtDx), 2026-11-13; not yet checked by eye.
\scanpage{24/19}
\mspage{19}

Indeed it suffices to verify this for 2.12 and this follows
from 2.11 and from the fact that for an algebraic prescheme in
characteristic zero $\rightarrow$ smooth $=$ regular.

\uline{Section 3} (or 20.3)*
\uline{Generic hyperplane section} : \uline{geometric irreducibility and}
\uline{connectedness.}
\uline{Theorem 3.1} (Bertini - Zariski)

Let us assume dim $f(X) \geq 2$ and that X is geometrically
irreducible. Then the generic hyperplane section $\underline{Y}_{\eta}$ has the same
property.

Let K/k be the function field of X and let n = dim $\underline{P}$;
introducing the affine coordinates $T_1,\dots,T_n$ in $\underline{P}$ (by choosing a
hyperplane at infinity H$\infty$ such that f(X) is not contained in it)
and $S_1,\dots,S_n$ the affine coordinates in $\underline{P}$v, we see that the
function field L of $Y_{\eta}$ can be identified with the field of
fractions of the integral domain $K[S_1,\dots,S_n] \,/\, (\Sigma\, t_iS_i - 1)$ where
the $t_i \in K$ are the images of $T_i$ with the aid of $f{:}X \rightarrow \underline{P}$. Since
dim $\overline{f(X)} \geq 0$, the $t_i$ are not all algebraic over k, á fortiori they
are not all zero let us assume for example that $t_n \neq 0$. We realize
then immediately that we have $L = K(S_1,\dots,S_n)$ (pure transcendental
extension), $S_n \in L$ being given by the equation $\Sigma\, t_iS_i - 1 = 0$
[en fonction Fr] as a function of the $S_i$ $(1 \leq i \leq n - 1)$ and the $t_i$
$(1 \leq\ i \leq n)$. On the other hand $k^1 = k(\eta)$ can be identified with
$k(S_1,\dots,S_n)$ and the canonical inclusion $k^1 \rightarrow L$ can be obtained by
sending $S_i$ to $S_i$ [PB: check this!] i.e. $k^1$ as a subextension of L

\margin{* Ask Grothendieck about the numbering system.}
